% Generated draft: requires mathematical and source review.
\documentclass[11pt]{article}
\usepackage[margin=25mm]{geometry}
\usepackage{fontspec}
\setmainfont{Latin Modern Roman}
\usepackage{amsmath,amssymb,mathtools,unicode-math}
\setmathfont{Latin Modern Math}
\usepackage{enumitem,needspace,xurl,hyperref}
\setlength{\parindent}{0pt}
\setlength{\parskip}{5pt}
\setlength{\emergencystretch}{3em}
\newcommand{\N}{\mathbb{N}}
\newcommand{\Z}{\mathbb{Z}}
\newcommand{\Q}{\mathbb{Q}}
\newcommand{\R}{\mathbb{R}}
\newcommand{\C}{\mathbb{C}}
\newcommand{\F}{\mathbb{F}}
\newcommand{\A}{\mathbb{A}}
\newcommand{\PP}{\mathbb{P}}
\newcommand{\E}{\mathbb{E}}
\newcommand{\eps}{\varepsilon}
\newcommand{\dd}{\,\mathrm{d}}
\begin{document}
\section*{UnsolvedMath ID: 30006649}
\textbf{Source problem code:} catalogue entry
\subsection{Definitions and mathematical statement}
Let \(p\) and \(q\) be distinct primes. A pro-\(p\) group is a topological group isomorphic to an inverse limit of finite \(p\)-groups. For such a group \(G\), let \(G^{\delta}\) denote its underlying abstract group, with the topology forgotten. Let \(\F_q\) be the field with \(q\) elements.

The supplied statement does not explicitly specify a \(G\)-action on the coefficients. The formulation below uses the usual interpretation of unadorned coefficients: \(\F_q\) is a trivial \(G\)-module. Confirmation of this interpretation from the source context is required. With this convention, define
\[
H^n_{\mathrm{abs}}(G;\F_q):=H^n(G^{\delta};\F_q),
\]
where the right-hand side is ordinary group cohomology, computed using all group cochains, without a continuity requirement. Write \(H^n_{\mathrm{cts}}(G;\F_q)\) for continuous group cohomology, using continuous cochains and giving \(\F_q\) the discrete topology.

\textbf{Question, under the trivial-action interpretation.} Is it true that for every pair of distinct primes \(p,q\), every pro-\(p\) group \(G\), and every integer \(n\geq 1\),
\[
H^n_{\mathrm{abs}}(G;\F_q)=0?
\]
The supplied record compares this proposed vanishing with the continuous-cohomology vanishing
\[
H^n_{\mathrm{cts}}(G;\F_q)=0\qquad(n\geq 1)
\]
under the same coefficient convention. No finite-generation or other additional restriction on \(G\) appears in the supplied reconstructed statement. The verbatim source fragment is incomplete, so its missing context cannot independently confirm these hypotheses.
\subsection{Short English statement}
For every pro-\(p\) group and every prime \(q\neq p\), must ordinary group cohomology with trivial \(\F_q\)-coefficients vanish in every positive degree when continuity is not imposed on cochains? The trivial coefficient action is implicit in the supplied notation and requires confirmation. The record labels the question open; this draft does not independently verify that status.
\subsection{Sources}
\begin{itemize}
\item Local UnsolvedMath registry, permanent ID \(30006649\), exact source problem code \texttt{catalogue entry}, in the supplied Oberwolfach Reports Open Problems collection. No individual proposer is identified in the record.
\item \emph{Cohomology of Finite Groups: Interactions and Applications} (2026), Oberwolfach Reports, pp.\ 541--610, DOI \texttt{10.4171/owr/2026/9}. Supplied originating report: \url{https://doi.org/10.4171/owr/2026/9}. The supplied verbatim fragment includes page number 551 but omits the antecedent of ``this'' and the surrounding definitions.
\item Ivanov--Mikhailov, \emph{On discrete homology of a free pro-p-group}: \url{https://arxiv.org/abs/1705.09131}. This is supplied as related literature, not as a resolution of the stated question.
\item Literature assessment dated exactly \texttt{2026-08-22}: the supplied curation labels the question open, says that the report records it as unresolved, and says that the cited discrete-cohomology work does not settle the general vanishing claim. This is supplied dated curation, not an independently verified literature assessment.
\item The registry labels the reconstructed statement \texttt{corrected\_verified} and states that it was reconstructed and checked against the official report. These are supplied provenance claims; this draft has not independently inspected either external source.
\end{itemize}
\end{document}
